\documentclass[11pt]{article}

\usepackage[T1]{fontenc}
\usepackage[utf8]{inputenc}
\usepackage{lmodern}
\usepackage[english]{babel}

\usepackage[margin=1in]{geometry}
\usepackage{amsmath,amssymb,amsthm}
\usepackage{booktabs}
\usepackage[expansion=false]{microtype}
\usepackage{enumitem}
\usepackage[colorlinks=true,allcolors=blue]{hyperref}
\setlength{\emergencystretch}{3em}

\newtheorem{theorem}{Theorem}
\newtheorem{lemma}{Lemma}
\newtheorem{corollary}{Corollary}
\newtheorem{proposition}{Proposition}
\newtheorem{assumption}{Assumption}
\theoremstyle{definition}
\newtheorem{definition}{Definition}
\theoremstyle{remark}
\newtheorem{remark}{Remark}

% --- shared notation with the base paper (Part I) ---
\newcommand{\lst}{\lambda_{\mathrm{st}}}
\newcommand{\lred}{\lambda_{\mathrm{red}}}
\newcommand{\lirr}{\lambda_{\mathrm{irr}}}
\newcommand{\Hraw}{H_{\mathrm{raw}}}
\newcommand{\Hgated}{H_{\mathrm{gated}}}
\newcommand{\am}{\bar\alpha_m}
\newcommand{\bm}{\bar\beta_m}
\newcommand{\eseg}{\varepsilon_{\mathrm{seg}}}
% --- new to this paper ---
\newcommand{\Esc}{E}
\newcommand{\fesc}{f}
\newcommand{\Stealth}{\mathrm{Stealth}}
\newcommand{\image}{\mathrm{im}}
\newcommand\blfootnote[1]{\begingroup\renewcommand\thefootnote{}\footnote{#1}\addtocounter{footnote}{-1}\endgroup}

\title{\vspace{-2em}Escape, Cost, and Correlation at the Verification Floor\\
\large of Gated Agentic Computation}

\author{Ivo Matija\v{s}evi\'c\\ \small Independent Researcher \\ \small ivomatijasevic@gmail.com}
\date{July 2026}

\begin{document}
\maketitle
\blfootnote{\textcopyright~2026 Ivo Matija\v{s}evi\'c.\ This work is licensed under a Creative Commons
Attribution 4.0 International (CC BY 4.0) license, \url{https://creativecommons.org/licenses/by/4.0/}.}

\begin{abstract}
Two companion papers establish a \emph{verification floor} for gated agentic systems: \emph{A
Fault-Tolerance Threshold for Gated Agentic Computation}~\cite{matijasevic2026n1} shows that a gate
family's shared blind spot---a stealth mass $\lst$, decomposing as $\lst=\lred+\lirr$---caps the
reliable horizon at $\Hgated\approx\Hraw/\lst$ regardless of gate count; and \emph{Generated Gates
Inherit Their Generator's Blind Spots}~\cite{matijasevic2026gg} shows that a system generating its own
gates cannot fall below that floor, which lowers only as \emph{exogenous} model families or
\emph{reality-grounded} gates are imported. Taking that floor as given, this paper develops three
consequences at the floor and draws one practical framing; the underlying identities and inequalities
are classical and attributed to their sources.
\textbf{(i)~A series-system escape identity.} Modelling an output as $n$ units each carrying an
undetected fault with (fixed) probability $\fesc$---the accept-conditioned per-unit escape
probability---the chance the output contains at least one escaped fault is
$\Esc(n)=1-(1-\fesc)^n\to1$ for any fixed $\fesc>0$; the size-axis reading of Part~I's horizon ceiling,
with the same stealth-mass floor on $\fesc$.
\textbf{(ii)~A conditional cost comparison.} Under an assumed power-law error--compute relation with
exponent $p$ and a per-family stealth multiplier, reaching a target escape rate a distance $\delta$
above the shared floor costs $O(\log(1/\delta))$ by gating versus $O((1/\delta)^{1/p})$ by scaling the
executor's own error down---a model calculation that formalizes the currently empirical ``verify rather
than scale'' claim.
\textbf{(iii)~A correlated-escape bound.} When output units are monotone functions of shared independent
upstream faults, escape events are positively associated, so $\Esc_{\mathrm{corr}}(n)\le 1-(1-\fesc)^n$
for any number of shared parents---an instance of the classical association / series-system bound, with
a mutual-independence equality characterization.
Finally, the floor explains a proposed ordinal assurance ladder (why internal rigor cannot substitute
for independent validation) and a typed-evidence non-collapse rule. We give the epistemic/diagnosability
reading of the floor that ties these to established theory, state each result's assumptions and honest
caveats, and flag open questions.
\end{abstract}

\section{Introduction}
Part~I~\cite{matijasevic2026n1} asked how far \emph{verification} can extend the reliable horizon of
a fixed, unreliable executor, and answered with three results in a standard constant-hazard model:
a logarithmic-overhead threshold theorem, a $1/\lst$ horizon ceiling set by the verification stack's
shared blind spot, and a Young--Daly cost-optimal checkpoint interval.
The headline of that work is that the binding constraint on autonomous horizon is not raw model
capability but the \emph{diversity} of the verification stack.

This paper takes that verification floor---established statistically in~\cite{matijasevic2026n1} and
mechanistically in the companion~\cite{matijasevic2026gg}---as given, and develops three of its
consequences plus a practical framing. None is claimed as new mathematics; the identities and
inequalities are classical, attributed in place and in \S\ref{sec:related}. What is ours is the
synthesis and, in \S\ref{sec:cost}, one closed-form cost comparison we did not find stated elsewhere.
Our contributions:
\begin{itemize}[leftmargin=1.4em,itemsep=2pt]
\item \textbf{A size-axis reading of the horizon ceiling (\S\ref{sec:escape}).} The standard
series-system escape identity $\Esc(n)=1-(1-\fesc)^n$, read as the output-size analogue of Part~I's
time-axis ceiling, making precise that ``zero faults at scale'' requires driving the per-unit floor to
zero, not adding verification.
\item \textbf{A conditional cost comparison (\S\ref{sec:cost}).} A model calculation comparing the cost
of reaching a target escape rate by gating versus by scaling the executor's own error, with the shared
floor made explicit; it formalizes a currently empirical folk claim.
\item \textbf{A correlated-escape bound (\S\ref{sec:corr}).} The classical association / series-system
bound applied to output provenance: shared provenance, at fixed unit marginals, can only lower the
probability of at least one escape relative to independence, for any number of shared parents, with a
mutual-independence equality characterization.
\item \textbf{An assurance framing the floor grounds (\S\ref{sec:assurance}).} A proposed ordinal
evidence ladder and a typed-evidence non-collapse rule---both light repackagings of assurance-case
practice---which the floor explains: internal rigor is bounded by it, and heterogeneous evidence
carries structure a scalar destroys.
\end{itemize}
Along the way, \S\ref{sec:epistemic} records an epistemic/diagnosability reading of the floor that
connects it to established theory (epistemic logic, rough sets, discrete-event diagnosability). We
recap Part~I's notation (\S\ref{sec:recap}), develop the three results
(\S\S\ref{sec:escape}--\ref{sec:corr}), draw the assurance framing (\S\ref{sec:assurance}), and close
with open problems (\S\ref{sec:open}) and related work (\S\ref{sec:related}).

\section{The model, recapped}\label{sec:recap}
We adopt the model and notation of Part~I~\cite{matijasevic2026n1} and restate it only to fix terms;
the reader is referred there for the three theorems and their proofs.

Agentic work proceeds at a constant fault hazard $\lambda$ per unit of work and is divided into
checkpointed \emph{segments} of length $s$; a segment is \emph{bad} with probability
$p=1-e^{-\lambda s}$.
At each checkpoint a stack of $m$ imperfect \emph{gates}, combined by majority vote, decides whether
to accept or recompute the segment; each gate has a false-accept rate $\alpha$ and a false-reject rate
$\beta$ (T1 below assumes both $<\tfrac12$), and $\am,\bm$ denote the effective majority-vote rates of
the stack. Because a rejected segment is recomputed, the emitted segment's error rate is conditioned on
acceptance---Part~I's per-segment slip probability $\eseg$, which we reuse in \S\ref{sec:escape}.
A \emph{gate family} is a set of gates sharing a mode of judgement.
The family's \emph{stealth mass} $\lst$ is the fraction of bad artifacts that the entire family
accepts---its shared blind spot; it decomposes as $\lst=\lirr+(1-\lirr)\rho$ into a
specification/intent-bound \emph{irreducible} floor $\lirr$ (faults that satisfy the checkable
contract but violate intent) and a lineage-correlated \emph{reducible} part governed by a per-family
probability $\rho$, with $k$ independent families driving the reducible part to $\rho^k\to0$ but
never touching $\lirr$.

Part~I established, in this model: (T1) a threshold theorem---\emph{when the stealth mass is zero}
($\lst=0$) and $\max(\alpha,\beta)<\tfrac12$, a stack of $m=O(\log(T/(\varepsilon s)))$ gates per
checkpoint achieves end-to-end success $\ge1-\varepsilon$ over horizon $T$ at multiplicative overhead
$O(\log T)$; (T2) a horizon ceiling---when $\lst>0$, $\Hgated\le\Hraw/\lst$ for every $m$, and
$\Hgated\le\Hraw/\lirr$ for any diversity breadth;
and (T3) a cost-optimal checkpoint interval $s^\star\approx\sqrt{G/\lambda}$ with $G=mg$ the
per-checkpoint verification cost.
We build on (T1)--(T3) throughout and do not reprove them.

\section{The series-system escape identity}\label{sec:escape}
Part~I bounds reliability along the \emph{time} axis: an output produced over a horizon, verified
segment by segment.
Many outputs are better described by their \emph{size}: a document of $n$ claims, a program of $n$
functions, a data record of $n$ fields.
We give the size-axis analogue and show it subsumes the horizon ceiling.

We first fix the sampling process, since the right per-unit quantity depends on it. A candidate unit
is generated and passed to the verification gate; on rejection it is recomputed and re-checked until a
candidate is accepted, and that accepted candidate is \emph{emitted}. We assume the retry attempts are
i.i.d.\ draws from the same candidate-and-gate law with $\Pr[\text{accept}]>0$, so the emitted unit is a
generic candidate conditioned on acceptance. Write $\pi$ for the prevalence of
a bad candidate, $a_B=\Pr[\text{accept}\mid\text{bad}]$ for the family's false-accept rate on bad
candidates, and $a_G=\Pr[\text{accept}\mid\text{good}]$ for its acceptance rate on good ones (its
false-reject rate is $1-a_G$). An \emph{emitted} unit is one that passed the gate, so its escape
probability is conditioned on acceptance: the relevant per-unit quantity is the posterior
\begin{equation}\label{eq:fdef}
\fesc\;:=\;\Pr[\text{bad}\mid\text{accepted}]\;=\;\frac{\pi\,a_B}{\pi\,a_B+(1-\pi)\,a_G}.
\end{equation}
This is exactly Part~I's per-segment slip probability $\eseg$ (accept-conditioned), read per output unit
rather than per checkpoint. Since $\fesc$ is non-decreasing in $a_B$ and a bad artifact in the stealth set
is accepted by every gate---hence by any monotone aggregator that accepts on unanimous acceptance
($A(1,\dots,1)=1$), majority included---the family's false-accept rate satisfies $a_B\ge\lst$, so for a
\emph{fixed} stack $\fesc$ is floored: with $0<\pi<1$ and $a_G\le1$,
\begin{equation}\label{eq:ffloor}
\fesc\;\ge\;\frac{\pi\lst}{\pi\lst+(1-\pi)a_G}\;\ge\;\frac{\pi\lst}{\pi\lst+(1-\pi)}\;>\;0\quad\text{whenever }\lst>0.
\end{equation}
The $a_G$-dependent middle bound is the sharp floor for a fixed stack; the right-hand ($a_G=1$) bound is
the stack-independent family floor. No amount of same-family verification lowers $a_B$ below $\lst$, so
none lowers $\fesc$ below the stack-independent floor. The levers that \emph{do} lower it act on $\lst$
(grounding, diversity, and specifications, \S\ref{sec:epistemic}), on the prevalence $\pi$ (a better
executor), or on $a_G$ (fewer false rejects, which dilute the accepted pool).

\begin{proposition}[Series-system escape identity]\label{thm:escape}
Model an output as $n$ emitted units, each independently carrying an undetected fault with fixed
probability $\fesc\in[0,1]$ from~\eqref{eq:fdef}. Then the probability that the output contains at least
one escaped fault is the series-system reliability complement
\begin{equation}\label{eq:escape}
\Esc(n)=1-(1-\fesc)^n\;=\;n\fesc-\binom{n}{2}\fesc^2+\cdots,
\end{equation}
with the union bound $\Esc(n)\le n\fesc$ (an inequality, tight to first order for $n\fesc\ll1$).
\end{proposition}
\begin{proof}
Each unit is clean with probability $1-\fesc$; by independence all $n$ are clean with probability
$(1-\fesc)^n$, and $\Esc(n)$ is the exact complement.
\end{proof}

Equation~\eqref{eq:escape} is entirely classical---the reliability of a series system of $n$
independent components~\cite{barlowproschan1965}, equivalently the independent family-wise error of
simultaneous inference (the \v{S}id\'ak form~\cite{sidak1967}) and the same at-least-one-event form as
the code model $\mathrm{pass}@k=1-(1-p)^k$~\cite{chen2021,brown2024} (with escapes in place of successes). It is also, read token-by-token, the
much-discussed autoregressive-divergence argument that $\Pr[\text{correct}]=(1-e)^n\to0$; a recent
critique~\cite{arbuzov2025} names this exact model ``the prevailing assumption'' and argues its
uniform-independence premise fails because errors concentrate at a sparse set of key tokens. We make
no novelty claim for~\eqref{eq:escape}; the observation we add is only its reading as the size-axis
form of Part~I's horizon ceiling.

\begin{corollary}[Zero-escape needs a zero floor]\label{cor:zero}
For \emph{fixed} $\fesc$, $\Esc(n)\to0$ as $n\to\infty$ iff $\fesc=0$; for any fixed $\fesc>0$,
$\Esc(n)\to1$. (If $\fesc=\fesc_n$ varies with $n$, the governing condition is $n\fesc_n\to0$.)
\end{corollary}

The corollary is the honest ceiling on ``no faults at scale.'' Because $\fesc$ increases in $a_B$ and
same-family gates keep $a_B\ge\lst$, once $a_B$ has saturated at the shared blind spot, adding more
gates of the \emph{same family} cannot lower $\fesc$ past the floor~\eqref{eq:ffloor}
(Part~I~\cite{matijasevic2026n1}). Diversity lowers the floor only through the \emph{reducible} part of
the stealth mass: $k$ independent families drive the stealth-limited floor down along
$\lst^{(1{:}k)}=\lirr+(1-\lirr)\rho^{\,k}\to\lirr$ (Part~I's decomposition, with the reducible residuals
independent across families conditional on the common irreducible set---distinct from the unconditional
exogenous-family product $\mu(\Sigma_F\cap\Sigma_N)=\lambda_F\lambda_N$ of~\cite{matijasevic2026gg}), so
the floor on $a_B$---and hence on
$\fesc$---approaches its $\lirr$-level value, \emph{not} zero. (Whether the operational $a_B$ itself
attains that floor depends on the aggregator; for the unanimous model $a_B=\lst$.) The irreducible floor $\lirr$ yields only to better specifications
or reality-grounding; the prevalence $\pi$ yields only to a better executor. Below saturation,
same-family depth does reduce the reducible residual---the claim is only that it cannot pass the floor.

\paragraph{Unification.}
Equation~\eqref{eq:escape} is one quantity into which separate results collapse. Because
$\fesc$~\eqref{eq:fdef} is precisely Part~I's per-segment slip probability $\eseg$, the horizon ceiling
of Part~I is the time-ordered instance of the same geometric law: an output is a sequence of $n=T/s$
segments, each emitted with slip probability $\fesc=\eseg$, and end-to-end success $(1-\fesc)^n$ is
\eqref{eq:escape} read along the time axis. The two are the same identity on two different indices, and
$\lst$ supplies the floor on $\fesc$ in both. A size-axis generalization to outputs whose units share
provenance is the subject of \S\ref{sec:corr}; a generalization to arbitrary provenance \emph{graphs}
is left open (\S\ref{sec:open}).

\paragraph{Evidence and caveats.}
The identity is illustrated in the companion script \texttt{simulate\_floor.py} at $\fesc=0.05$:
at $n=20$ the predicted escape $0.6415$ matches the Monte-Carlo value $0.6407$ (seed $20260718$) within
sampling error, and the fit holds across $n\in\{1,5,10,20,50\}$. Two modelling assumptions are worth
naming. First, \eqref{eq:escape} assumes the $n$ units are \emph{independent}; units passing a common
terminal gate share that gate's blind spot, a correlation channel distinct from the provenance sharing
of \S\ref{sec:corr}---and, \emph{to the extent that channel fits the monotone independent-primitive
model of Assumption~\ref{ass:prov}}, it pushes $\Esc$ below the independent value. Second, ``output
unit'' has a clean meaning for structured
outputs (fields, claims, functions), but its right definition for continuous or unstructured outputs is
open, and $\Esc(n)$ should be read with that caveat where units are not crisply individuated.

\section{The stealth-mass floor and its epistemic reading}\label{sec:epistemic}
The floor these results rest on is not established here; it is the subject of two companion papers.
Part~I~\cite{matijasevic2026n1} gives the statistical floor: a gate family's \emph{stealth mass}
$\lst$---the fraction of bad artifacts every gate in the family accepts---bounds majority accuracy, and
decomposes as $\lst=\lred+\lirr$ into a diversity-reducible part and a specification-bound irreducible
floor. The companion paper~\cite{matijasevic2026gg} gives the mechanism: modelling a family $F$ by its
representation relation $\sim_F$ (indistinguishability under every prompt), the stealth set
$\Sigma_F=\{\text{bad }c:\ c\sim_F c'\text{ for some good }c'\}$---with $\mu$ the fault-generating
distribution and $\lambda_F=\mu(\Sigma_F)$---is accepted by every \emph{sound} gate built within $F$, and
with nonzero false rejects up to a vanishing margin (its GG1). Consequently within-family generation
cannot drive the escape floor below $\lambda_F$---however deep the generation, however many gates check
gates, however diverse the prompts (its GG2; we use only this lower-bound direction, which is what our
conclusions need, and not GG2's achievability, which requires sufficient within-family coverage). Only
two moves lower the floor: importing an \emph{exogenous} family, whose blind spot, when independent,
multiplies it ($\mu(\Sigma_F\cap\Sigma_N)=\lambda_F\lambda_N$, driving $\lred$ down), and a
\emph{reality-grounded} gate---one deciding by ground truth rather than any family's
representation---which removes the grounded part of $\lirr$. We take this floor and its two levers as
given.

\paragraph{An epistemic/diagnosability reading.} The one observation we add is that this floor is a
familiar object across several established literatures---worth recording because it locates the
phenomenon and supplies external tools. The relation $\sim_F$ is an indistinguishability relation in the
possible-worlds sense; the accept set $\{w:[w]_F\cap G\neq\emptyset\}$ (the bad worlds $\sim_F$-linked to
some good one, whose $B$-mass is $\lambda_F$) is exactly the rough-set upper approximation of the good
set~\cite{pawlak1982}; ``a fault is undetectable iff observationally indistinguishable from acceptable
behaviour, and duplicating an observation channel without new information cannot restore detection'' is
the substance of diagnosability and decentralized fault-diagnosis for discrete-event
systems~\cite{sampath1995,debouk2000,descodiag2024}; and the pooled vision of diverse families is the
\emph{distributed knowledge} of multi-agent epistemic logic~\cite{fagin1995,halpernmoses1990} (the
intersection $\bigcap_i\sim_i$ of the family relations, strictly finer than any one). On the
language-model side the same impossibility has just appeared: Mason and Anand~\cite{masonanand2026} prove
that a text-only supervisor cannot distinguish grounded from fabricated outputs and that no amount of
same-channel supervision resolves it---GG2's within-family invariance in a different formalism---%
consistent with the empirical ``models think alike'' correlation-floor findings
of~\cite{goel2025,kim2025}. Every one of these readings names the same lever: a distinction must be
\emph{imported}---an exogenous family, or reality-grounding---never bred within a family. This is the
semantic content behind the assurance framing of \S\ref{sec:assurance} (``who ran the check, on whose
object'' is the informal name for ``which family's $\sim_F$ was applied'').

\section{A conditional cost comparison: gating versus scaling}\label{sec:cost}
Why gate an unreliable executor rather than wait for, or pay for, a more reliable one? This section
gives a model calculation---not a general theorem---comparing the cost of the two routes. It
formalizes a claim now argued empirically across the test-time-compute
literature~\cite{snell2024,lifshitz2025,scalingverif2026}: that adding verification is a cheaper lever
than scaling the generator. Everything here is conditional on the two cost models we state as
assumptions, and neither ingredient is individually new---logarithmic redundancy for a reliability
target is von Neumann's~\cite{vonneumann1956} and Part~I's, and the other side is inverting an
empirical power law.

We work throughout in the per-emitted-unit escape probability $\fesc$ of \S\ref{sec:escape}, and compare
the two routes \emph{for a fixed baseline executor}---so the prevalence $\pi$ is held fixed and the
irreducible floor is a constant. In the epistemic model that floor is the posterior value
$q_{\mathrm{irr}}=\pi\lirr/(\pi\lirr+(1-\pi)a_G)>0$, but we take $q_{\mathrm{irr}}$ as an exogenous
constant of a fixed specification and executor (on the scaling route, a genuinely better executor lowers
$\pi$ and hence $\fesc$ separately---a different lever, out of scope here). We compare the cost of
reaching a target a distance
\[
\delta\;:=\;q-q_{\mathrm{irr}}\;>\;0
\]
above the floor, where $q$ is the required per-unit escape; for an output of $N$ independent units a
target output-escape $\varepsilon$ sets $q\approx\varepsilon/N$ (\S\ref{sec:escape}). We compare as
$\delta\to0^+$.

\begin{assumption}[Scaling cost, floor-aware]\label{ass:det}
Driving the executor's own per-unit escape down by compute obeys a power law above the floor,
$e(C)=q_{\mathrm{irr}}+\kappa\,C^{-p}$ for constants $\kappa>0$, $p\in(0,1]$. We take the exponent from
language-model \emph{loss} scaling~\cite{kaplan2020,hoffmann2022} ($p\approx0.05$--$0.5$) as a
\emph{hypothesis only}: those laws fit cross-entropy loss, not task-error or intent-violation rate, and
the loss-to-error map is an extra unstated step.
\end{assumption}

\begin{assumption}[Gating cost, floor-aware]\label{ass:gate}
Diverse gate families reduce the \emph{reducible} conditional false-accept rate geometrically:
$a_B(n)=\lirr+(a_{B,0}-\lirr)\,c^{\,n}$ for a per-family multiplier $c\in(0,1)$ (Part~I's decomposition
with conditional independence across families; \S\ref{sec:escape}). We further assume the good-acceptance
rate is bounded away from zero, $a_G(n)\ge a_G^{\min}>0$, and $0<\pi<1$, $\fesc_0>q_{\mathrm{irr}}$. Under
these, the posterior map $a_B\mapsto\fesc$~\eqref{eq:fdef}---a M\"obius transform, smooth and increasing
with derivative bounded away from $0$ and $\infty$ on the range---carries the geometric decay to the
induced escape at the \emph{same rate},
\[
f_{\mathrm{gate}}(n)-q_{\mathrm{irr}}\;=\;\Theta(c^{\,n}),
\]
which is all Proposition~\ref{prop:cost} uses. We write $f_{\mathrm{gate}}(n)=q_{\mathrm{irr}}+(\fesc_0-q_{\mathrm{irr}})c^{\,n}$
below only as a representative closed form with this rate, not an exact identity; each family costs a
fixed $g_F$. The $a_G(n)\ge a_G^{\min}$ hypothesis is load-bearing: if instead conjoining families drives
$a_G(n)=\prod_i a_{G,i}\to0$ (pure unanimous stacking), then for $\lirr>0$ the posterior $\fesc\to1$, not
$q_{\mathrm{irr}}$---the comparison then fails, and holding $a_G$ up requires a calibrated aggregator, not
raw conjunction (see ``where the argument stops'').
\end{assumption}

\emph{Scale.} Solving $e(C)=q_{\mathrm{irr}}+\delta$ under Assumption~\ref{ass:det} gives the compute to make the
executor itself that reliable,
\begin{equation}\label{eq:det}
C(\delta)=\big(\kappa/\delta\big)^{1/p},
\end{equation}
polynomial in $1/\delta$ with a large exponent $1/p\approx2$--$20$.

\emph{Gate.} Solving $f_{\mathrm{gate}}(n)\le q_{\mathrm{irr}}+\delta$ gives the number of families,
\begin{equation}\label{eq:gate}
n_{\mathrm{gate}}(\delta)=\max\!\Big\{0,\ \Big\lceil\tfrac{\ln((\fesc_0-q_{\mathrm{irr}})/\delta)}{\ln(1/c)}\Big\rceil\Big\}
=O\!\big(\log(1/\delta)\big),
\end{equation}
logarithmic in $1/\delta$, at total cost $g_F\,n_{\mathrm{gate}}(\delta)$.

\begin{proposition}[Cost separation]\label{prop:cost}
Under Assumptions~\ref{ass:det}--\ref{ass:gate}, as the target approaches the shared floor
($\delta\to0^+$) the gating cost is logarithmic and the scaling cost polynomial in $1/\delta$, so their
ratio vanishes:
\[
\frac{g_F\,n_{\mathrm{gate}}(\delta)}{C(\delta)}
=\frac{g_F\,O(\log(1/\delta))}{(\kappa/\delta)^{1/p}}\;\longrightarrow\;0 .
\]
\end{proposition}
\begin{proof}
Put $x=1/\delta\to\infty$. Then $C=\kappa^{1/p}x^{1/p}$ and
$g_F\,n_{\mathrm{gate}}\le g_F\big(1+\ln((\fesc_0-q_{\mathrm{irr}})x)/\ln(1/c)\big)$, so the ratio is
$O(\log x)/x^{1/p}\to0$, since $\ln x=o(x^{1/p})$ for every $p>0$.
\end{proof}

The separation is meaningful for $\delta\in(0,\ \fesc_0-q_{\mathrm{irr}}]$; it says nothing at or below
the floor, where $\delta\le0$ and neither route helps. Table~\ref{tab:sep} makes the gap concrete for
one illustrative choice of constants ($p=0.1$; $\kappa=0.1$; $\fesc_0-q_{\mathrm{irr}}=0.1$;
$c=\tfrac12$; $g_F=1$);
the asymptotic separation, not the row values, is the result, and every entry is reproducible
from~\eqref{eq:det}--\eqref{eq:gate}.

\begin{table}[h]
\centering
\small
\begin{tabular}{lrr}
\toprule
excess target $\delta=q-q_{\mathrm{irr}}$ & scale compute $C(\delta)$ & gate families $n_{\mathrm{gate}}(\delta)$ \\
\midrule
$10^{-2}$ & $10^{10}$ & $4$ \\
$10^{-4}$ & $10^{30}$ & $10$ \\
$10^{-6}$ & $10^{50}$ & $17$ \\
$10^{-8}$ & $10^{70}$ & $24$ \\
\bottomrule
\end{tabular}
\caption{Illustrative cost of reaching an excess target $\delta$ above the shared floor by scaling the
executor ($C=(\kappa/\delta)^{1/p}$) versus by gating
($n_{\mathrm{gate}}=\lceil\log_c(\delta/(\fesc_0-q_{\mathrm{irr}}))\rceil$), for $p=0.1$, $\kappa=0.1$,
$\fesc_0-q_{\mathrm{irr}}=0.1$, $c=\tfrac12$. Scaling compute grows by ten orders of magnitude per decade of
$\delta$; gating grows by about three families.}\label{tab:sep}
\end{table}

\paragraph{Where the argument stops, and what it rests on.}
This is a \emph{toy asymptotic specialization}, not an operational cost model, and it instantiates the
classical redundancy-versus-component-improvement trade-off of reliability engineering~\cite{kuo2000}
in the Part~I floor metric; we do not claim the trade-off itself, only its statement here. Four
assumptions are load-bearing. First, the power law of Assumption~\ref{ass:det} is a hypothesis carried
over from loss scaling, not a proven law for task error. Second, $q_{\mathrm{irr}}$ is assumed exogenous
and common to both routes; on the scaling route, lowering the executor's error changes the prevalence
$\pi$, so a fixed floor there is itself an assumption (Corollary~\ref{cor:zero}'s ``only specs/grounding
lower the floor'' holds at a fixed executor; a better executor lowers $\pi$, hence $\fesc$, separately).
Third, the comparison counts only executor compute $C$ against gate cost $g_F n$; it omits the
\emph{retry cost} of gating---conjoining $n$ families multiplies the good-acceptance rate to
$\prod_i a_{G,i}$, so the expected attempts per emitted unit, $1/\!\Pr[\text{accept}]$, can grow with
$n$ and erode the advantage unless false rejects are controlled (e.g.\ by a majority rather than
unanimous aggregator, which then breaks the clean product law)---as well as baseline executor compute,
diversity-acquisition cost, and latency. Fourth, Assumption~\ref{ass:gate} presumes an \emph{unbounded
supply} of independent families, exactly the scarce resource the rest of this theory names as binding.
Finally, we mention but do \emph{not} rely on the automata metaphor sometimes attached to this
comparison: determinizing a nondeterministic automaton by subset construction blows up
exponentially~\cite{rabinscott1959}, and the blow-up is genuinely necessary~\cite{meyerfischer1971}.
That is a statement about representational state complexity, not about paying compute to lower a
stochastic error rate; it is at most an illustrative analogy and is not evidence for
Proposition~\ref{prop:cost}.

\section{Correlated escape under shared provenance}\label{sec:corr}
The escape identity (\S\ref{sec:escape}) assumes independent output units. Real outputs often share
upstream provenance: many fields derived from one computation, many claims resting on one retrieved
fact. A shared upstream fault is then a \emph{single} event, not $n$ independent ones, and one expects
correlated escape to sit \emph{below} the independent bound. This is a textbook consequence of the
classical theory of \emph{association} of random variables~\cite{esary1967} and of coherent-system
reliability~\cite{barlowproschan1965}: monotone functions of independent variables are associated, and
associated components satisfy exactly the series-system inequality below. We do not claim the
inequality; our contribution is only its application to output provenance and the packaged strictness
characterization, which the companion papers state only for the independent-unit case.

\begin{assumption}[Monotone provenance model]\label{ass:prov}
The escape of unit $j$ is an indicator $I_j=\phi_j(X)\in\{0,1\}$, where
$X=(X_1,\dots,X_M)\in\{0,1\}^M$ is a vector of \emph{mutually independent} primitive fault/leak
variables---the faults of the shared upstream parents together with each unit's own local fault
source, each surviving its gates independently---and each $\phi_j$ is coordinatewise non-decreasing.
Monotonicity encodes that a primitive fault can only ever cause, never prevent, an escape: $I_j=1$
if any parent feeding unit $j$ leaks past its upstream gates, or its local source does. Write
$\fesc_j=\Pr[I_j=1]$, and $\fesc$ for the common value when the units are exchangeable.
\end{assumption}

The bound is the association / series-system inequality~\cite{esary1967,barlowproschan1965}; we state
the product form used here and prove it for completeness (its cleanest source is~\cite{esary1967}).

\begin{lemma}[Association product form; Esary--Proschan--Walkup]\label{lem:harris}
Let $X=(X_1,\dots,X_M)$ have independent coordinates, and let $g_1,\dots,g_n:\{0,1\}^M\to[0,\infty)$
be bounded and all coordinatewise non-increasing (equivalently, by symmetry, all non-decreasing). Then
$\mathbb{E}\big[\prod_{j=1}^n g_j(X)\big]\ge\prod_{j=1}^n\mathbb{E}[g_j(X)]$.
\end{lemma}
\begin{proof}
The base case $n=2$ is Harris's inequality~\cite{harris1960} (a special case of FKG~\cite{fkg1971}):
for independent coordinates, two bounded functions monotone in the same direction satisfy
$\mathbb{E}[g_1 g_2]\ge\mathbb{E}[g_1]\mathbb{E}[g_2]$. The non-decreasing and non-increasing cases are
interchanged by replacing each $X_i$ with $1-X_i$, which preserves independence and flips every $g_j$
simultaneously; we argue the non-increasing case. For the induction, the partial product
$G=\prod_{j<n}g_j$ is again non-negative, bounded, and non-increasing (a product of non-negative
non-increasing functions), so the $n=2$ case applied to $(G,g_n)$ and the induction hypothesis give
$\mathbb{E}[\prod_{j\le n}g_j]=\mathbb{E}[G\,g_n]\ge\mathbb{E}[G]\,\mathbb{E}[g_n]
\ge\big(\prod_{j<n}\mathbb{E}[g_j]\big)\mathbb{E}[g_n]$. This is the closure of association under
coordinatewise-monotone maps and the series-system bound of~\cite{esary1967}.
\end{proof}

\begin{proposition}[Correlated-escape bound]\label{prop:corr}
Under Assumption~\ref{ass:prov}, for \emph{any} number of shared parents,
\begin{equation}\label{eq:corr}
\Esc_{\mathrm{corr}}(n)\;=\;1-\Pr\!\Big[\textstyle\bigcap_{j=1}^{n}\{I_j=0\}\Big]
\;\le\;1-\prod_{j=1}^{n}(1-\fesc_j),
\end{equation}
which equals $1-(1-\fesc)^n$ when the units share a common marginal $\fesc$.
Assume in addition every unit has $\fesc_j<1$. Then the bound is strict as soon as some shared parent
$X_s$ is \emph{jointly pivotal} for two distinct units $j,k$---i.e.\ $X_s$ is non-degenerate and, on a
positive-probability set of the other coordinates $X_{-s}$, both $\phi_j$ and $\phi_k$ are non-constant
in $X_s$---and $n\ge2$. Under the same $\fesc_j<1$ nondegeneracy, equality holds \emph{iff} the $I_j$
are mutually independent.
\end{proposition}
\begin{proof}
Each $g_j:=1-I_j=1-\phi_j(X)$ is a bounded, non-negative, coordinatewise non-increasing function of
the independent vector $X$. Lemma~\ref{lem:harris} gives
\[
\Pr\!\Big[\textstyle\bigcap_{j}\{I_j=0\}\Big]
=\mathbb{E}\Big[\textstyle\prod_{j}(1-I_j)\Big]
\;\ge\;\prod_{j}\mathbb{E}[1-I_j]=\prod_{j}(1-\fesc_j),
\]
and taking complements yields~\eqref{eq:corr}; with exchangeable units $\prod_j(1-\fesc_j)=(1-\fesc)^n$.
This uses only monotonicity and independence of the \emph{primitives}, so it holds for any number of
shared parents.

\emph{Strictness.} Let $U=1-I_j$, $V=1-I_k$. By the law of total covariance,
\[
\operatorname{Cov}(U,V)=\underbrace{\mathbb{E}\big[\operatorname{Cov}(U,V\mid X_{-s})\big]}_{(\mathrm I)}
+\underbrace{\operatorname{Cov}\big(\mathbb{E}[U\mid X_{-s}],\ \mathbb{E}[V\mid X_{-s}]\big)}_{(\mathrm{II})}.
\]
Term $(\mathrm{II})\ge0$: $\mathbb{E}[U\mid X_{-s}]$ and $\mathbb{E}[V\mid X_{-s}]$ are non-increasing
functions of the independent coordinates $X_{-s}$, so Harris applies. Term $(\mathrm{I})>0$: under
joint pivotality, on a positive-probability set of $X_{-s}$ both $U$ and $V$ are non-constant and
same-direction monotone in the non-degenerate $X_s$, giving strictly positive conditional covariance
there and $\ge0$ elsewhere. Hence $\operatorname{Cov}(U,V)>0$, i.e.\
$\mathbb{E}[UV]>(1-\fesc_j)(1-\fesc_k)$. Treating $h:=UV$ as one non-negative non-increasing factor and
applying Lemma~\ref{lem:harris} to $\{h\}\cup\{1-I_l:l\neq j,k\}$,
\[
\mathbb{E}\Big[\textstyle\prod_l(1-I_l)\Big]\;\ge\;\mathbb{E}[h]\!\!\prod_{l\neq j,k}\!\!\mathbb{E}[1-I_l]
\;>\;(1-\fesc_j)(1-\fesc_k)\!\!\prod_{l\neq j,k}\!\!(1-\fesc_l)=\prod_l(1-\fesc_l),
\]
where the strict step uses $\fesc_l<1$ so that $\prod_{l\neq j,k}(1-\fesc_l)>0$. Thus~\eqref{eq:corr} is
strict. \emph{Equality:} mutual independence gives $\mathbb{E}[\prod_j(1-I_j)]=\prod_j(1-\fesc_j)$ at
once. Conversely, suppose equality holds and all $\fesc_l<1$. For \emph{any} subset $S\subseteq\{1,\dots,n\}$,
group $h_S=\prod_{j\in S}(1-I_j)$ (non-negative and non-increasing) and apply Lemma~\ref{lem:harris} to
$\{h_S\}\cup\{1-I_l:l\notin S\}$:
\[
\prod_l(1-\fesc_l)=\mathbb{E}\Big[\textstyle\prod_l(1-I_l)\Big]\;\ge\;\mathbb{E}[h_S]\!\!\prod_{l\notin S}\!(1-\fesc_l)\;\ge\;\prod_l(1-\fesc_l),
\]
so both inequalities are equalities; since $\prod_{l\notin S}(1-\fesc_l)>0$ (all $\fesc_l<1$), this
forces $\mathbb{E}[h_S]=\prod_{j\in S}(1-\fesc_j)$, i.e.\ $\Pr[\bigcap_{j\in S}\{I_j=0\}]=\prod_{j\in S}(1-\fesc_j)$.
Holding for every $S$, this is mutual independence of the events $\{I_j=0\}$, hence of the $I_j$. (We do
not need, or invoke, any pairwise-implies-mutual result for associated variables; the subset form of
Lemma~\ref{lem:harris} delivers mutual independence directly.)
\end{proof}

\begin{remark}[The multi-parent case]
One might ask whether ``strictly below independent'' survives when a unit has $k\ge3$ shared parents,
where multi-parent dependence might turn an escape correlation negative. It cannot: the escape
indicators are monotone functions of \emph{independent} primitives, so association holds for \emph{any}
$k$~\cite{esary1967}, and strictness is governed only by joint pivotality (with $\fesc_l<1$), not by
parent count. This was never open in the classical association literature---EPW places no restriction on
the number of shared inputs; we simply make it explicit here for the provenance setting.
\end{remark}

\begin{remark}[Interpretation]\label{rem:interp}
Equation~\eqref{eq:corr} compares dependence structures \emph{at fixed unit marginals}: positive
common-cause dependence \emph{clusters} escapes, so the probability of at least one escape is no larger
than under independence. It does \emph{not} say shared provenance ``helps reliability.'' Shared causes
raise the chance of simultaneous multi-unit failure, destroy redundancy benefits, and in practice often
\emph{raise} the marginals themselves---the harmful regime studied under common-cause failure. The
bound is a statement about the OR-event under fixed marginals, nothing more.
\end{remark}

\paragraph{Evidence and scope.}
The bound is illustrated in \texttt{simulate\_floor.py} on a one-shared-parent event model (a shared
upstream fault present with $q_{\mathrm{up}}=0.3$, surviving the upstream gate with $a_{\mathrm{up}}=0.247$
and each output's sink with $a_{\mathrm{sink}}=0.247$; independent local faults with $q_{\mathrm{loc}}=0.05$,
per-unit marginal $\fesc_{\mathrm{marg}}=0.0304$). At $n=20$ the correlated escape is $0.2777$ (exact;
Monte-Carlo $0.2785$, seed $20260718$) against the independent bound $0.4610$---a gap of $0.183$ that
widens across the displayed range (both probabilities tend to $1$ as $n\to\infty$, so the gap
eventually narrows). The \emph{problem} of where to place verification on a provenance graph is already
public~\cite{ro2025sherlock}; we stop at the bound and leave any allocation policy out of scope
(\S\ref{sec:open}).

\section{Grading assurance: a framing the floor explains}\label{sec:assurance}
The results here rest on one object---an escape floor that internal effort cannot cross---and that
floor has a practical corollary for how one should \emph{talk} about whether a gated system
``works.'' We propose, not prove, an organizing frame with two parts, both deliberately modest and both
already anticipated in the assurance literature; what this paper adds is the observation that the floor
results explain \emph{why} each holds.

\subsection{A proposed evidence ladder}
Claims that a built system ``works'' or is ``better than before'' can be graded on an ordinal ladder by
\emph{who ran the check, on whose object, for whose benefit}: (R1)~\textbf{built and internally
verified}---demonstrated by a re-runnable check the builder wrote, on an object the builder controls;
(R2)~\textbf{independently reproduced}---a second party reproduces the result from the builder's
artifact without the builder's help; (R3)~\textbf{externally validated}---exercised by a party outside
the builder's control, on a workload the builder did not author; (R4)~\textbf{value-validated}---that
external party derives measured value from it, not merely confirms it runs. We offer this as a
\emph{proposed} taxonomy, not an exhaustive law, and it is a coarsening of frameworks that grade along
related independence and maturity axes: independent verification and validation (IV\&V, graded by
evaluator independence), the reproduced/replicated distinction of artifact-evaluation badging (graded by
who ran the check, on whose artifacts), technology-readiness levels~\cite{mankins1995} (maturity and
operating environment), typed assurance-case evidence~\cite{gsn}, evidence-based levels of evidence,
standardized AI assurance
levels~\cite{aal2026}, and the evidence-licensing framework of Li~\cite{li2026}. We claim no new
taxonomy; the axes---evaluator independence, workload independence, and demonstrated utility---are in
fact partly orthogonal, so the rungs should be read as a rough, claim-relative progression rather than
a total order (a weak R4 signal is not automatically stronger than a rigorous R1 proof). Its only
virtue is a cheap test---the three questions above---for the common overclaim of quoting internal rigor
as external assurance.

The floor results say why the rungs cannot collapse into internal rigor, though the link is a heuristic
proxy, not a theorem. A system whose internal checks all share one blind structure can have arbitrarily
many of them and still sit at R1: same-family depth cannot lower the escape floor
(\S\ref{sec:escape}, Cor.~\ref{cor:zero}; the self-generation floor of~\cite{matijasevic2026gg}). Moving outward is
an \emph{opportunity} to add a genuinely new indistinguishability relation---an external validator
(R3) or paying user (R4) exercises the system under assumptions and workloads the builder's gates never
saw---but organizational independence does not \emph{guarantee} epistemic non-redundancy: an R2
reproducer re-running the builder's own check largely re-applies the builder's relation and mainly
hardens the \emph{evidence} (against fraud or setup error), and an external party can share the
builder's assumptions and add nothing. Actual floor reduction requires an imported family $N$ whose
stealth set $\Sigma_N$ genuinely excludes part of the existing blind spot, $\mu(\Sigma_F\setminus\Sigma_N)>0$
(\S\ref{sec:epistemic}); the ladder is a proxy for seeking one, not a proof that a higher rung supplies it.

\subsection{Typed evidence should not collapse to a scalar}
A pipeline that mixes a conditional formal proof, a deterministic test or compile result, a
probabilistic model-judgment (an LLM or human reviewer's opinion), and a provenance claim should not
report a single external-facing scalar---a probability or ``trust score''---folding all four together.
This is standard in safety-case practice, where confidence is argued from typed, unmerged
evidence rather than reduced to one attribute~\cite{gsn,bloomfield2022}, and it is being restated for
language-model pipelines as typed evidence graphs~\cite{wang2026}. (Evidence-combination formalisms
such as Dempster--Shafer~\cite{shafer1976} do produce numbers, but they retain ignorance and conflict
rather than collapsing to a single score---an illustration that a defensible aggregate must keep its
semantics, not that no number may ever appear.) The rule is thus not ``never quote a number'' but:
any external aggregate must carry explicit semantics and stay linked to the typed evidence, provenance,
and assumptions it came from; a bare scalar that hides evidence type is what must not be surfaced.

The epistemic model of \S\ref{sec:epistemic} offers a useful interpretation. Different evidence types
are different indistinguishability relations with different blind spots; a \emph{bare} scalar fold is a
projection that does not ordinarily expose which part of the confidence is load-bearing under which
conditions---the structure the distributed-knowledge reading of \S\ref{sec:epistemic} retains.
A conditional theorem and a measured reviewer rate are sound under different
envelopes; averaging them into an unlabelled number reports neither honestly. Non-collapse is the reporting-side image of the
paper's recurring fact that diversity lives in the \emph{structure} of what different gates can and
cannot see, not in a scalar. Neither the ladder nor this rule is new---each is a light repackaging of
assurance-case and evidence-theory practice~\cite{mankins1995,gsn,shafer1976}; the contribution is only
the link to the floor.

\section{Open problems}\label{sec:open}
Five questions are adjacent to the results above but not settled here; we state them as open, with no
result claimed, to mark the intended shape of the theory.

\begin{enumerate}[leftmargin=1.6em,itemsep=3pt]
\item \textbf{A provenance-graph escape law.} The escape identity (\S\ref{sec:escape}) treats
independent units and \S\ref{sec:corr} treats units that are monotone Boolean functions of arbitrary
\emph{independent} primitive faults (any number of shared parents, one layer). The genuinely open case
is a full provenance DAG with \emph{internal} gates, statistically \emph{dependent} primitives,
recomputation/retries, and non-monotone repairs, and a sink-floor law relating escape to graph
structure and per-node floors. We regard this as a prove-first item and do not include it here.
\item \textbf{An exact finite-horizon escape process.} Equations~\eqref{eq:escape}
and~\eqref{eq:corr} are static bounds; the exact escape probability of a checkpointed process with
recomputation is naturally a finite-horizon Markov decision problem whose optimal gating policy would
refine the asymptotics of \S\ref{sec:cost}. This is in progress and deferred.
\item \textbf{Reviewer--producer capability matching.} The results here treat the executor and its
gates as given; how gate capability should be matched to executor capability (when a weaker but
diverse reviewer helps, when it does not) is open.
\item \textbf{A telemetry estimator for the stealth mass.} Part~I notes that the reducible stealth
mass is partly estimable from a system's own logs and targeted audits; a principled estimator for
$\lst$ from telemetry, with an identifiability analysis in the latent-class tradition of Hui and
Walter~\cite{huiwalter1980} (estimating error rates without a gold standard), is left to future work.
This is the quantity that would make the stealth mass predictable rather than only observable.
\item \textbf{Checker strength.} The framing of \S\ref{sec:assurance} treats gates as given. A separate
question is how the \emph{strength of the checking procedure}---whether the property it checks is
decidable now, corroborable soon, or gradeable only later---bounds what ``verified'' can honestly mean
for that check's output, and whether it predicts when a checkable-notation effort finds a consumer. We
regard this as a single idea deserving its own careful treatment and do not develop it here.
\end{enumerate}

\section{Related work}\label{sec:related}
This paper builds on Part~I~\cite{matijasevic2026n1} and its companion~\cite{matijasevic2026gg} (the
statistical and mechanistic accounts of the floor, respectively), and inherits their lineage: reliable
computation from unreliable components~\cite{vonneumann1956,moore1956}, the quantum threshold
theorem~\cite{aharonov1997}, Condorcet's jury theorem~\cite{condorcet1785}, the Young--Daly
checkpoint interval~\cite{young1974,daly2006}, and the $N$-version-programming tradition on coincident
failure~\cite{eckhardt1985,littlewood1989,knight1986}, whose contemporary agentic echo is the
large-scale $N$-version coding-agent study of Ron, Baudry and Monperrus~\cite{ron2026}.

For the specific results here, we attribute the prior art plainly, since the mathematics in each is
classical. \S\ref{sec:escape}: the escape identity~\eqref{eq:escape} is series-system
reliability~\cite{barlowproschan1965}, the \v{S}id\'ak independent-error form~\cite{sidak1967}, and the
at-least-one form of $\mathrm{pass}@k$~\cite{chen2021,brown2024}; token-wise it is the autoregressive-%
divergence argument whose independence premise~\cite{arbuzov2025} disputes. The correlation floor it
feeds is actively studied---Turkmen, Buyukates and Bastopcu~\cite{turkmen2026} derive an
information-theoretic ensemble error floor under a Gaussian-copula model, and ``models think alike''
correlation is measured by~\cite{goel2025,kim2025}. \S\ref{sec:epistemic}: the stealth-set floor, its
inheritance and self-generation invariance, and its two levers are established in the
companion~\cite{matijasevic2026gg} (statistically in Part~I); we only record its reading in
possible-worlds epistemic logic~\cite{fagin1995,halpernmoses1990}, rough-set upper
approximation~\cite{pawlak1982}, and discrete-event diagnosability~\cite{sampath1995,debouk2000,descodiag2024},
noting the language-model impossibility of Mason--Anand~\cite{masonanand2026}. \S\ref{sec:cost}: the comparison formalizes a claim
argued empirically across the test-time-compute literature~\cite{snell2024,lifshitz2025,scalingverif2026},
resting on loss scaling~\cite{kaplan2020,hoffmann2022}; the automaton-determinization
blow-up~\cite{rabinscott1959,meyerfischer1971} is a metaphor, not evidence. \S\ref{sec:corr}: the
bound is the classical association / series-system inequality~\cite{esary1967,harris1960,fkg1971,barlowproschan1965};
the provenance-placement \emph{problem} is treated by Sherlock~\cite{ro2025sherlock}; ours is only the
application to output provenance and the packaged strictness characterization.

\section{Conclusion}
The classical tools here converge on one object, the verification floor of Part~I
and~\cite{matijasevic2026gg}. At fixed prevalence,
same-family verification cannot lower a gate family's escape floor (\S\ref{sec:escape}; the floor is
\emph{set by} Part~I's stealth mass, and is not a consequence of the series identity itself, which
merely propagates it across output size); a system's own self-generated gates cannot lower the shared
blind spot (the self-generation floor of~\cite{matijasevic2026gg}); above that floor, under stated and admittedly strong cost models, gating
reaches a target more cheaply than scaling the executor's own error does (Prop.~\ref{prop:cost}); and
at fixed marginals, positive provenance dependence only clusters escapes, never raising the chance of
at least one relative to independence (Prop.~\ref{prop:corr}). None of these is new mathematics; the
contribution is the synthesis---three classical tools trained on the floor of Part~I
and~\cite{matijasevic2026gg}---plus the one conditional cost calculation. The floor has several levers:
diversity lowers the lineage-\emph{reducible} stealth mass, better specifications or reality-grounding
lower the irreducible part $\lirr$, and a stronger executor lowers the prevalence $\pi$; ``irreducible''
is always relative to a fixed executor
and specification. Much of this is verification engineering available today, alongside gains in raw
model capability.

\section*{Code and data availability}
The numerical claims of \S\S\ref{sec:escape}--\ref{sec:corr} are reproduced by
\texttt{simulate\_floor.py} at \url{https://github.com/ivmat/gated-computation-sim}. This paper builds
on two companion papers archived on Zenodo: Part~I,
DOI~\href{https://doi.org/10.5281/zenodo.20820968}{10.5281/zenodo.20820968}, and the self-generation
floor paper~\cite{matijasevic2026gg},
DOI~\href{https://doi.org/10.5281/zenodo.20837102}{10.5281/zenodo.20837102}.

\begin{thebibliography}{99}
\small
\bibitem{matijasevic2026gg} I.~Matija\v{s}evi\'c. Generated gates inherit their generator's blind
spots: a diversity floor for self-generated verification. Zenodo, 2026.
DOI:~10.5281/zenodo.20837102.
\bibitem{matijasevic2026n1} I.~Matija\v{s}evi\'c. A Fault-Tolerance Threshold for Gated Agentic
Computation: Reliable Long-Horizon Work from Unreliable Executors. 2026.
Archived on Zenodo, DOI:~10.5281/zenodo.20820968.
\url{https://github.com/ivmat/gated-computation-sim}.
\bibitem{vonneumann1956} J.~von Neumann. Probabilistic logics and the synthesis of reliable
organisms from unreliable components. In \emph{Automata Studies}, Princeton Univ.\ Press, 1956.
\bibitem{moore1956} E.~F.~Moore and C.~E.~Shannon. Reliable circuits using less reliable relays.
\emph{J.~Franklin Inst.}, 1956.
\bibitem{aharonov1997} D.~Aharonov and M.~Ben-Or. Fault-tolerant quantum computation with constant
error. In \emph{STOC}, 1997.
\bibitem{condorcet1785} M.~de Condorcet. \emph{Essai sur l'application de l'analyse \`a la
probabilit\'e des d\'ecisions rendues \`a la pluralit\'e des voix.} 1785.
\bibitem{young1974} J.~W.~Young. A first order approximation to the optimum checkpoint interval.
\emph{Comm.\ ACM}, 1974.
\bibitem{daly2006} J.~T.~Daly. A higher order estimate of the optimum checkpoint interval for
restart dumps. \emph{Future Generation Computer Systems}, 2006.
\bibitem{eckhardt1985} D.~E.~Eckhardt and L.~D.~Lee. A theoretical basis for the analysis of
multiversion software subject to coincident errors. \emph{IEEE Trans.\ Software Eng.}, 1985.
\bibitem{littlewood1989} B.~Littlewood and D.~Miller. Conceptual modeling of coincident failures
in multiversion software. \emph{IEEE Trans.\ Software Eng.}, 1989.
\bibitem{knight1986} J.~C.~Knight and N.~G.~Leveson. An experimental evaluation of the assumption
of independence in multiversion programming. \emph{IEEE Trans.\ Software Eng.}, 1986.
\bibitem{ron2026} J.~Ron, B.~Baudry, and M.~Monperrus. N-Version Programming with Coding Agents.
\emph{arXiv:2606.20158}, 2026.
\bibitem{fagin1995} R.~Fagin, J.~Y.~Halpern, Y.~Moses, and M.~Y.~Vardi. \emph{Reasoning About
Knowledge.} MIT Press, 1995.
\bibitem{kaplan2020} J.~Kaplan, S.~McCandlish, T.~Henighan, et al. Scaling Laws for Neural Language
Models. \emph{arXiv:2001.08361}, 2020.
\bibitem{hoffmann2022} J.~Hoffmann, S.~Borgeaud, A.~Mensch, et al. Training Compute-Optimal Large
Language Models. \emph{arXiv:2203.15556}, 2022.
\bibitem{rabinscott1959} M.~O.~Rabin and D.~Scott. Finite automata and their decision problems.
\emph{IBM J.\ Res.\ Develop.}, 1959.
\bibitem{turkmen2026} Y.~Turkmen, B.~Buyukates, and M.~Bastopcu. Don't Always Pick the
Highest-Performing Model: An Information Theoretic View of LLM Ensemble Selection.
\emph{arXiv:2602.08003}, 2026.
\bibitem{ro2025sherlock} Y.~Ro, H.~Qiu, \'I.~Goiri, R.~Fonseca, R.~Bianchini, A.~Akella, Z.~Wang,
M.~Erez, and E.~Choukse. Sherlock: Reliable and Efficient Agentic Workflow Execution.
\emph{arXiv:2511.00330}, 2025.
\bibitem{harris1960} T.~E.~Harris. A lower bound for the critical probability in a certain
percolation process. \emph{Proc.\ Cambridge Philos.\ Soc.}, 56(1):13--20, 1960.
\bibitem{fkg1971} C.~M.~Fortuin, P.~W.~Kasteleyn, and J.~Ginibre. Correlation inequalities on some
partially ordered sets. \emph{Comm.\ Math.\ Phys.}, 1971.
\bibitem{huiwalter1980} S.~L.~Hui and S.~D.~Walter. Estimating the error rates of diagnostic tests.
\emph{Biometrics}, 1980.
% --- reliability / association ---
\bibitem{barlowproschan1965} R.~E.~Barlow and F.~Proschan. \emph{Mathematical Theory of Reliability.}
Wiley, 1965.
\bibitem{esary1967} J.~D.~Esary, F.~Proschan, and D.~W.~Walkup. Association of random variables, with
applications. \emph{Ann.\ Math.\ Statist.}, 38(5):1466--1474, 1967.
\bibitem{sidak1967} Z.~\v{S}id\'ak. Rectangular confidence regions for the means of multivariate
normal distributions. \emph{J.~Amer.\ Statist.\ Assoc.}, 62(318):626--633, 1967.
% --- epistemic logic / diagnosability / rough sets ---
\bibitem{halpernmoses1990} J.~Y.~Halpern and Y.~Moses. Knowledge and common knowledge in a distributed
environment. \emph{J.~ACM}, 37(3):549--587, 1990.
\bibitem{pawlak1982} Z.~Pawlak. Rough sets. \emph{Int.\ J.\ Computer \& Information Sciences},
11:341--356, 1982.
\bibitem{sampath1995} M.~Sampath, R.~Sengupta, S.~Lafortune, K.~Sinnamohideen, and D.~Teneketzis.
Diagnosability of discrete-event systems. \emph{IEEE Trans.\ Automatic Control}, 40(9):1555--1575, 1995.
\bibitem{debouk2000} R.~Debouk, S.~Lafortune, and D.~Teneketzis. Coordinated decentralized protocols
for failure diagnosis of discrete event systems. \emph{Discrete Event Dynamic Systems}, 10:33--86, 2000.
\bibitem{descodiag2024} B.~Cui, Z.~Ma, S.~Li, and X.~Yin. On epistemic properties in discrete-event
systems: a uniform framework and its applications. \emph{arXiv:2409.06588}, 2024.
% --- automata ---
\bibitem{meyerfischer1971} A.~R.~Meyer and M.~J.~Fischer. Economy of description by automata, grammars,
and formal systems. In \emph{12th Annual Symposium on Switching and Automata Theory (SWAT/FOCS)}, 1971.
% --- scaling / test-time compute / LLM error correlation ---
\bibitem{chen2021} M.~Chen et al. Evaluating large language models trained on code.
\emph{arXiv:2107.03374}, 2021.
\bibitem{brown2024} B.~Brown, J.~Juravsky, R.~Ehrlich, R.~Clark, Q.~V.~Le, C.~R\'e, and A.~Mirhoseini.
Large Language Monkeys: scaling inference compute with repeated sampling. \emph{arXiv:2407.21787}, 2024.
\bibitem{snell2024} C.~Snell, J.~Lee, K.~Xu, and A.~Kumar. Scaling LLM test-time compute optimally can
be more effective than scaling model parameters. \emph{arXiv:2408.03314}, 2024.
\bibitem{lifshitz2025} S.~Lifshitz, S.~McIlraith, and Y.~Du. Multi-agent verification: scaling test-time
compute with multiple verifiers. \emph{arXiv:2502.20379}, 2025.
\bibitem{scalingverif2026} J.~Kwok, X.~Zhang, M.~Xu, Y.~Liu, A.~Mirhoseini, C.~Finn, and M.~Pavone.
Scaling verification can be more effective than scaling policy learning for vision--language--action
alignment. \emph{arXiv:2602.12281}, 2026.
\bibitem{arbuzov2025} M.~L.~Arbuzov, S.~Bei, Z.~Dong, D.~Kalaev, and A.~A.~Shvets. Beyond exponential
decay: rethinking error accumulation in large language models. \emph{arXiv:2505.24187}, 2025.
\bibitem{goel2025} S.~Goel, J.~Str\"uber, I.~A.~Auzina, et al. Great models think alike and this
undermines AI oversight. \emph{arXiv:2502.04313}, 2025.
\bibitem{kim2025} E.~Kim, A.~Garg, K.~Peng, and N.~Garg. Correlated errors in large language models.
\emph{arXiv:2506.07962}, 2025.
\bibitem{masonanand2026} T.~Mason and V.~Anand. Epistemic observability in language models.
\emph{arXiv:2603.20531}, 2026.
% --- assurance / evidence ---
\bibitem{mankins1995} J.~C.~Mankins. Technology readiness levels: a white paper. NASA Office of Space
Access and Technology, 1995.
\bibitem{gsn} T.~Kelly and R.~Weaver. The Goal Structuring Notation---a safety argument notation. In
\emph{Proc.\ DSN Workshop on Assurance Cases}, 2004. (See also the GSN Community Standard, v3, 2021.)
\bibitem{shafer1976} G.~Shafer. \emph{A Mathematical Theory of Evidence.} Princeton Univ.\ Press, 1976.
\bibitem{li2026} H.~Li. The calibration turn in AI-assisted research: a conceptual and methodological
framework for evidence-licensed claims. \emph{arXiv:2606.31273}, 2026.
\bibitem{wang2026} Y.~Wang et al. From agent traces to trust: a survey of evidence tracing and
execution provenance in LLM agents. \emph{arXiv:2606.04990}, 2026.
\bibitem{bloomfield2022} R.~Bloomfield and J.~Rushby. Assessing confidence with Assurance 2.0.
\emph{arXiv:2205.04522}, 2022.
\bibitem{aal2026} M.~Brundage et al. Frontier AI auditing: toward rigorous third-party assessment of
safety and security practices at leading AI companies. \emph{arXiv:2601.11699}, 2026.
\bibitem{kuo2000} W.~Kuo and V.~R.~Prasad. An annotated overview of system-reliability optimization.
\emph{IEEE Trans.\ Reliability}, 49(2):176--187, 2000.
\end{thebibliography}

\end{document}
